\documentclass[12pt,oneside]{book}

% ==========================================================
% METADATOS
% ==========================================================
\newcommand{\universidad}{Universidad de Buenos Aires (UBA)}
\newcommand{\facultad}{Facultad de Derecho}
\newcommand{\carrera}{Derecho Civil --- Parte General}
\newcommand{\materia}{Derechos Reales}
\newcommand{\titulo}{Teoría General: Adquisición, Transmisión y Extinción de los Derechos Reales}
\newcommand{\autor}{Isaac Edgar Camacho Ocampo}
\newcommand{\anio}{2020}

% ==========================================================
% PAQUETES
% ==========================================================
\usepackage[utf8]{inputenc}
\usepackage[T1]{fontenc}
\usepackage[spanish,es-nodecimaldot]{babel}

\usepackage[
    a4paper,
    top=2.5cm,
    bottom=2.5cm,
    left=3cm,
    right=2.5cm,
    headheight=15pt
]{geometry}

\usepackage{amsmath}
\usepackage{amssymb}
\usepackage{mathtools}

\usepackage{graphicx}
\usepackage{float}
\usepackage{caption}
\usepackage{subcaption}

\usepackage{booktabs}
\usepackage{array}
\usepackage{longtable}
\usepackage{tabularx}

\usepackage{enumitem}

\usepackage{xcolor}
\usepackage[most]{tcolorbox}

\usepackage{fancyhdr}
\usepackage{ragged2e}

\usepackage[
    colorlinks=true,
    linkcolor=blue!50!black,
    citecolor=green!40!black,
    urlcolor=blue!60!black,
    pdftitle={Teoría General: Adquisición, Transmisión y Extinción de los Derechos Reales},
    pdfauthor={Isaac Edgar Camacho Ocampo},
    pdfsubject={Apuntes universitarios},
    bookmarks=true
]{hyperref}

% ==========================================================
% CONFIGURACIÓN
% ==========================================================
\graphicspath{
    {./img/}
    {./graficos/}
    {./figuras/}
}

\setcounter{tocdepth}{2}

\setlength{\parindent}{0pt}
\setlength{\parskip}{0.5em}

\setlist[itemize]{
    topsep=0.3em,
    itemsep=0.2em
}

\setlist[enumerate]{
    topsep=0.3em,
    itemsep=0.2em
}

\clubpenalty=10000
\widowpenalty=10000

\pagestyle{fancy}
\fancyhf{}
\fancyhead[L]{\nouppercase{\leftmark}}
\fancyhead[R]{\materia}
\fancyfoot[C]{\thepage}
\renewcommand{\headrulewidth}{0.4pt}
\renewcommand{\footrulewidth}{0pt}

\fancypagestyle{plain}{
    \fancyhf{}
    \fancyfoot[C]{\thepage}
    \renewcommand{\headrulewidth}{0pt}
}

% ==========================================================
% COMANDOS
% ==========================================================
\newcommand{\abs}[1]{\left|#1\right|}
\newcommand{\norm}[1]{\left\lVert#1\right\rVert}
\newcommand{\vect}[1]{\mathbf{#1}}
\newcommand{\deriv}[2]{\frac{d#1}{d#2}}
\newcommand{\pderiv}[2]{\frac{\partial#1}{\partial#2}}

% ==========================================================
% ENTORNOS / CAJAS
% ==========================================================
\newtcolorbox{definition}[1]{
    enhanced,
    breakable,
    title={Definición: #1},
    fonttitle=\bfseries,
    colback=blue!3,
    colframe=blue!50!black,
    boxrule=0.8pt,
    arc=2mm
}

\newtcolorbox{important}{
    enhanced,
    breakable,
    title={Importante},
    fonttitle=\bfseries,
    colback=yellow!8,
    colframe=orange!70!black,
    boxrule=0.8pt,
    arc=2mm
}

\newtcolorbox{example}[1]{
    enhanced,
    breakable,
    title={Ejemplo: #1},
    fonttitle=\bfseries,
    colback=green!4,
    colframe=green!40!black,
    boxrule=0.8pt,
    arc=2mm
}

\newtcolorbox{formula}[1]{
    enhanced,
    breakable,
    title={Fórmula / Regla: #1},
    fonttitle=\bfseries,
    colback=purple!4,
    colframe=purple!50!black,
    boxrule=0.8pt,
    arc=2mm
}

\newtcolorbox{exercise}[1]{
    enhanced,
    breakable,
    title={Ejercicio: #1},
    fonttitle=\bfseries,
    colback=gray!5,
    colframe=gray!60!black,
    boxrule=0.8pt,
    arc=2mm
}

\newtcolorbox{note}{
    enhanced,
    breakable,
    title={Nota},
    fonttitle=\bfseries,
    colback=white,
    colframe=gray!50,
    boxrule=0.6pt,
    arc=2mm
}

% ==========================================================
% DOCUMENTO
% ==========================================================
\begin{document}

% ---------------------------------------------------------
% PORTADA
% ---------------------------------------------------------
\begin{titlepage}

\thispagestyle{empty}

\begin{center}

\vspace*{1cm}

{\large \textsc{\universidad}}

\vspace{0.2cm}

{\large \textsc{\facultad}}

\vspace{0.2cm}

{\large \carrera}

\vspace{3cm}

{\Huge \textbf{\materia}}

\vspace{1cm}

{\LARGE \titulo}

\vspace{0.8cm}

\textit{Comprender cada concepto es el primer paso para hacerlo propio.}

\vspace{1.2cm}

\rule{0.70\textwidth}{0.5pt}

\vspace{0.5cm}

{\large Apuntes de estudio}

\vspace{0.5cm}

\rule{0.70\textwidth}{0.5pt}

\vfill

{\large \autor}

\vspace{0.3cm}

{\large \anio}

\end{center}

\vfill

\begin{flushleft}
\textit{Este es un modesto aporte a los estudiantes de ``Derechos Reales'' no pretende sustituir el material oficial de la materia.}
\end{flushleft}

\end{titlepage}

% ---------------------------------------------------------
% ÍNDICE
% ---------------------------------------------------------
\pagenumbering{roman}
\tableofcontents
\clearpage
\pagenumbering{arabic}

% ==========================================================
% CAPÍTULO 1
% ==========================================================
\chapter[El sistema de los derechos reales]{El sistema de los derechos reales}

\section{Numerus clausus y ley federal}

Los derechos reales se distinguen con nitidez de los derechos personales, tanto por su naturaleza como por los caracteres propios de cada esfera. Constituyen un vínculo directo e inmediato entre un titular y una cosa. Su número es limitado: los derechos reales son exactamente \textbf{14 y solo 14} (principio de \textit{numerus clausus}).

A diferencia de los derechos personales, los derechos reales son \textbf{creados exclusivamente por la ley} y hoy se encuentran todos codificados. Se trata de la \textbf{ley federal}: en materia de derechos reales solamente la Nación es competente para legislar.

\section{Normas de orden público y sistema de propiedad}

Los particulares deben ceñirse estrictamente a la regulación legal de cada derecho real, porque las normas que los rigen son de \textbf{orden público}.

El fundamento de este carácter es que el derecho real de cualquier país responde a un esquema de sociedad único para toda la nación. Tiene que ver con el modo en que una nación organiza su sistema de propiedad y su sistema de defensa de las relaciones de poder respecto de las cosas; con el modo en que una sociedad define el alcance de las facultades que está dispuesta a reconocer a cada titular del derecho real; y con las limitaciones que está dispuesta a aceptar en función de los intereses de la comunidad.

\begin{important}
\textbf{Artículo 12 (CCyCN).} Las convenciones particulares no pueden dejar sin efecto las leyes en cuya observancia está interesado el orden público. Es el principio de orden público aplicado a los derechos reales.
\end{important}

\section{Constitucionalización e internacionalización del derecho privado}

El Código Civil y Comercial de la Nación (CCyCN) evidencia un verdadero cambio de paradigma respecto de la legislación de derechos privados anterior a su vigencia, esencialmente porque incorpora los tratados internacionales de derechos humanos como ley vigente.

\textbf{Artículos 1 y 2.} Son fundamentos del Código y, a la vez, una regla de interpretación: los tratados de derechos humanos integran el ordenamiento y todo debe ser mirado bajo esa perspectiva.

\textbf{Artículo 15.} Las personas son titulares de los derechos individuales sobre todos los bienes que integran su patrimonio. Este artículo expresa el sistema de propiedad privada como eje del ordenamiento.

\textbf{Artículo 14.} La ley no ampara el ejercicio abusivo de los derechos individuales cuando puede afectar al ambiente y a los derechos de incidencia colectiva. Es un límite al ejercicio de los derechos individuales.

\textbf{Artículo 240.} Establece de manera rotunda un límite a los derechos individuales cuando resulta necesario compatibilizarlos con los derechos de incidencia colectiva: el derecho al ambiente, el derecho a respirar. Consagra la preeminencia del interés colectivo.

\begin{important}
En el sistema jurídico argentino existe una clara \textbf{preeminencia del interés colectivo por sobre el interés individual}. Por eso, en ocasiones resulta necesario modelar el alcance del ejercicio de un derecho real en función de los intereses de la comunidad.
\end{important}

\section{Cuadro Cornell: el sistema de los derechos reales}

% TODO: en las notas originales el cuadro Cornell afirma que los tratados se incorporan "con igual jerarquía constitucional"; el desarrollo de la clase solo indica que se los incorpora como "ley vigente". Se conservó la formulación del desarrollo.

\begin{longtable}{>{\raggedright\arraybackslash}p{0.27\textwidth}>{\raggedright\arraybackslash}p{0.65\textwidth}}
\toprule
\textbf{Preguntas / palabras clave} & \textbf{Notas / respuestas} \\
\midrule
\endfirsthead
\toprule
\textbf{Preguntas / palabras clave} & \textbf{Notas / respuestas} \\
\midrule
\endhead
\bottomrule
\endfoot
\textbf{¿Qué son los derechos reales y cuántos existen?} &
Son derechos subjetivos que establecen una relación directa e inmediata entre un titular y una cosa. Existen exactamente 14 (\textit{numerus clausus}), creados exclusivamente por la ley federal. Las normas que los regulan son de orden público y no pueden ser modificadas por los particulares. \\[0.8em]

\textbf{¿Por qué los derechos reales son de orden público?} &
Porque responden al esquema de organización de la propiedad de toda la nación. El Estado define el alcance de las facultades reconocidas al titular y las limitaciones en función del interés colectivo. Las convenciones privadas no pueden apartarse de esa regulación (Art.~12). \\[0.8em]

\textbf{¿Qué es la constitucionalización del derecho privado?} &
Es la incorporación de los tratados internacionales de derechos humanos como ley vigente. Los Arts.~1 y 2 del CCyCN los establecen como fundamento y regla de interpretación. Los Arts.~14 y 240 ponen límites al ejercicio individual de los derechos cuando afecta al ambiente o a los derechos de incidencia colectiva. \\
\midrule
\multicolumn{2}{p{\dimexpr0.92\textwidth+2\tabcolsep\relax}}{%
\textbf{Resumen:} los derechos reales son 14, creados solo por la ley federal y regidos por normas de orden público (Art.~12), porque expresan el modo en que la nación organiza su sistema de propiedad. El Código incorpora los tratados de derechos humanos (Arts.~1 y 2) y consagra la preeminencia del interés colectivo sobre el individual (Arts.~14 y 240).} \\
\end{longtable}

% ==========================================================
% CAPÍTULO 2
% ==========================================================
\chapter[Adquisición de los derechos reales]{Adquisición de los derechos reales}

\section{Adquisición originaria y derivada}

El sistema reconoce dos formas de adquisición de los derechos reales.

\begin{table}[H]
\centering
\small
\begin{tabularx}{\textwidth}{>{\raggedright\arraybackslash}p{0.22\textwidth}>{\raggedright\arraybackslash}X>{\raggedright\arraybackslash}p{0.26\textwidth}}
\toprule
\textbf{Tipo} & \textbf{Descripción} & \textbf{Ejemplo} \\
\midrule
Adquisición originaria &
El sistema reconoce la posibilidad de adquirir un derecho real en forma libre y absoluta, sin que derive de un vínculo anterior ni de ningún vínculo jurídico previo. &
Cosas sin dueño (\textit{res nullius}); prescripción adquisitiva. \\[0.5em]
Adquisición derivada (entre vivos) &
El derecho real es transmitido por un enajenante a un adquirente: uno transmite y el otro adquiere. &
Compraventa de un inmueble. \\[0.5em]
Adquisición derivada (causa de muerte) &
La transmisión opera por el fallecimiento del titular y sus bienes pasan al heredero. &
Sucesión hereditaria. \\
\bottomrule
\end{tabularx}
\caption{Formas de adquisición de los derechos reales}
\end{table}

\textbf{Artículo 1896.} El juez no puede constituir derechos reales ni imponer su constitución. Es una expresión del principio de legalidad en la constitución de los derechos reales: el derecho real se adquiere por los modos previstos en la ley, no por decisión judicial.

\section{Regla \textit{nemo plus iuris}}

\begin{formula}{Artículo 399 --- \textit{nemo plus iuris}}
Nadie puede transmitir a otro un derecho mejor o más extenso que el que tiene.
\end{formula}

La regla \textit{nemo plus iuris} es una herramienta fundamental del sistema. Se aplica tanto a los derechos reales como a los derechos personales y, si bien tiene excepciones previstas en la ley, constituye la regla general.

\section{Teoría del título y modo}

La adquisición derivada de los derechos reales por actos entre vivos requiere necesariamente de dos elementos, lo que la doctrina denomina \textbf{teoría del título y modo}.

\begin{important}
\textbf{Artículo 1892.} El derecho real se adquiere con la concurrencia del \textbf{título suficiente} más el \textbf{modo suficiente}. Es el fundamento central de la adquisición derivada de derechos reales por actos entre vivos.
\end{important}

\subsection{El título suficiente}

El título no debe identificarse en su sentido instrumental (el instrumento en papel), sino en su sentido abstracto e intangible.

\begin{definition}{Título suficiente}
Es el acto jurídico revestido de las condiciones de fondo y de forma establecidas por la ley, que tiene por finalidad transmitir o constituir un derecho real. El acto jurídico causal de la transmisión constituye el título suficiente. Esa causa puede ser la compraventa, la cesión, la permuta o la dación en pago.
\end{definition}

\textbf{Artículo 259.} El acto jurídico es el acto voluntario lícito que tiene por fin inmediato la adquisición, modificación o extinción de las relaciones o situaciones jurídicas.

\textbf{Artículo 260.} El acto voluntario es el ejecutado con discernimiento, intención y libertad, que se manifiesta por un hecho exterior.

\textbf{Artículo 1017.} Para los inmuebles, el acto jurídico causal debe plasmarse necesariamente en \textbf{escritura pública}. Es la forma requerida para el título suficiente en materia de inmuebles.

\begin{note}
Para ilustrar que el título suficiente no es el instrumento en papel, se recurrió a una analogía: entregar un instrumento que no plasma un acto jurídico causal equivale a no tener título suficiente. El título suficiente es el acto jurídico intangible que plasma la causa de la transmisión.
\end{note}

\subsection{El modo suficiente: la tradición}

El modo suficiente es el elemento propio de los derechos reales que se ejercen por la posesión.

\begin{definition}{Tradición}
Es la entrega material de la cosa: una persona entrega la cosa a otra que la recibe. El poder de hecho que tenía el sujeto se transmite a quien la adquiere con la tradición. Es el modo suficiente en los derechos reales que se ejercen por la posesión.
\end{definition}

El \textbf{tradens} es quien entrega y el \textbf{accipiens} es quien recibe. La tradición es \textbf{constitutiva} en todos los casos en que se adquieren derechos reales que se ejercen por la posesión: sin tradición no se adquirió el derecho real.

\begin{important}
\textbf{Artículo 1924.} La tradición no puede suplirse por fórmulas escritas. Cuando en una escritura pública se establece que <<en este acto las partes se entregan la tradición>>, se estaría suplantando un acto material por una fórmula escrita en el papel, lo cual no puede hacerse. No puede transmitirse la posesión mediante palabras escritas.
\end{important}

\subsection{Inscripción registral: efectos declarativo y constitutivo}

El efecto de la inscripción registral depende del tipo de bien.

\begin{table}[H]
\centering
\small
\begin{tabularx}{\textwidth}{>{\raggedright\arraybackslash}p{0.17\textwidth}>{\raggedright\arraybackslash}p{0.2\textwidth}>{\raggedright\arraybackslash}X}
\toprule
\textbf{Bien} & \textbf{Efecto de la inscripción} & \textbf{Consecuencia práctica} \\
\midrule
Inmuebles & Declarativo &
El derecho nace con el título más el modo (escritura más tradición). La inscripción solo lo hace oponible \textit{erga omnes} (conocido por toda la comunidad). \\[0.5em]
Automotores & Constitutivo &
Sin la inscripción (transferencia registral) no se es dueño del automotor. El derecho nace con la inscripción. \\
\bottomrule
\end{tabularx}
\caption{Efectos de la inscripción registral}
\end{table}

\section{Cuadro Cornell: adquisición de los derechos reales}

\begin{longtable}{>{\raggedright\arraybackslash}p{0.27\textwidth}>{\raggedright\arraybackslash}p{0.65\textwidth}}
\toprule
\textbf{Preguntas / palabras clave} & \textbf{Notas / respuestas} \\
\midrule
\endfirsthead
\toprule
\textbf{Preguntas / palabras clave} & \textbf{Notas / respuestas} \\
\midrule
\endhead
\bottomrule
\endfoot
\textbf{¿Cuáles son los tipos de adquisición de los derechos reales?} &
Originaria: sin vínculo jurídico anterior (ej.: cosas sin dueño, prescripción). Derivada: por transmisión de un titular anterior; puede ser por actos entre vivos (compraventa) o por causa de muerte (sucesión hereditaria). \\[0.8em]

\textbf{¿Qué establece la regla \textit{nemo plus iuris} (Art.~399)?} &
Nadie puede transmitir a otro un derecho mejor o más extenso que el que tiene. Si el transmitente no era propietario o su derecho estaba limitado, el adquirente recibirá solo lo que aquel tenía, con las mismas limitaciones. \\[0.8em]

\textbf{¿Qué exige el Art.~1892 para adquirir un derecho real por actos entre vivos?} &
La concurrencia de título suficiente y modo suficiente (teoría del título y modo). \\[0.8em]

\textbf{¿Qué es el título suficiente?} &
Es el acto jurídico (en sentido abstracto, no el instrumento) revestido de las condiciones de fondo y forma exigidas por la ley, que tiene por finalidad transmitir o constituir un derecho real. En inmuebles debe plasmarse en escritura pública (Art.~1017). \\[0.8em]

\textbf{¿Qué es el modo suficiente?} &
Es la tradición: la entrega material efectiva de la cosa de una persona a otra. Es constitutiva en los derechos reales que se ejercen por la posesión. Sin tradición no se adquiere el derecho real, y no puede suplirse por fórmulas escritas (Art.~1924). \\[0.8em]

\textbf{¿Cuál es el efecto de la inscripción registral en inmuebles y automotores?} &
En inmuebles: efecto declarativo. El derecho nace con el título más el modo; la inscripción solo lo hace oponible a terceros. En automotores: efecto constitutivo. Sin inscripción (transferencia) no se adquiere el dominio. \\
\midrule
\multicolumn{2}{p{\dimexpr0.92\textwidth+2\tabcolsep\relax}}{%
\textbf{Resumen:} la adquisición puede ser originaria o derivada. La derivada entre vivos exige título suficiente (acto jurídico causal) y modo suficiente (tradición) según el Art.~1892, bajo la regla \textit{nemo plus iuris} (Art.~399). La tradición no se suple con fórmulas escritas (Art.~1924). La inscripción es declarativa en inmuebles y constitutiva en automotores.} \\
\end{longtable}

% ==========================================================
% CAPÍTULO 3
% ==========================================================
\chapter[Casos especiales de tradición]{Casos especiales de tradición}

Los dos únicos casos en los que el ordenamiento admite que la adquisición del derecho real opere sin la entrega material efectiva de la cosa son la \textbf{traditio brevi manu} y el \textbf{constituto posesorio}. Fuera de estos supuestos, previstos expresamente por la ley, el principio de legalidad de los derechos reales no admite apartarse del requisito de tradición.

\section{Traditio brevi manu}

\begin{definition}{Traditio brevi manu (Arts.~1892 y 1933)}
Es la excepción a la exigencia de tradición material que permite que quien tenía la cosa en virtud de otro título (por ejemplo, el locatario) pase a tenerla como propietario, sin necesidad de tradición material, haciendo constar el cambio en la escritura.
\end{definition}

\begin{example}{La inquilina que compra su departamento}
María es una mujer joven, docente y soltera, que durante muchísimos años alquiló un departamento en el barrio de Florida. En un momento se le presentó la posibilidad de comprar ese mismo departamento, lo cual coincidía con el interés del propietario en venderlo. María firma la escritura pública de compraventa (título suficiente), pero la ley exige también el modo: la tradición, es decir, la entrega material de la cosa.

\textbf{Sin la traditio brevi manu}, María debería llamar a la mudadora, sacar todos sus muebles, salir en calidad de inquilina, hacer que el propietario ingresara y que luego se lo entregara, cumpliendo los requisitos de propietario, capacidad y entrega. Sería un dispendio inútil.

\textbf{Solución legal:} María, que ya tenía la cosa a título de inquilina (tenencia), pasa a tenerla ahora como propietaria (dominio), haciendo constar este cambio en la escritura, sin necesidad de ningún movimiento material. Asciende de categoría: la flecha va hacia arriba ($\uparrow$).
\end{example}

\subsection{Variante: el nuevo adquirente y la inquilina que continúa}

% TODO: contenido incompleto o ambiguo en las notas originales. El apunte titula esta variante como "constituto posesorio" pero la desarrolla con la figura de la traditio brevi manu; se conserva el texto tal como figura.

Supóngase que no es María quien puede comprar, sino otro adquirente, y que María necesita seguir alquilando el departamento. En ese caso, el anterior propietario vende y el nuevo adquirente también debería recibir la posesión material. Sin embargo, María podría continuar en el inmueble como nueva inquilina del adquirente. La traditio brevi manu permite que María, que tenía la cosa a nombre del propietario original, pase a tenerla a nombre del adquirente sin movimiento material, con debida constancia notarial y fecha cierta.

\section{Constituto posesorio}

\begin{definition}{Constituto posesorio (Arts.~1892 y 1923)}
Es la segunda excepción a la exigencia de tradición material y el caso inverso a la traditio brevi manu: permite que quien poseía la cosa como propietario pase a tenerla como tenedor (por ejemplo, como nuevo inquilino), sin necesidad de tradición material y con la debida constancia notarial.
\end{definition}

\begin{example}{El propietario endeudado de la crisis de 2001}
En el año 2001, un señor con muchísimas deudas decide vender su departamento para pagarlas, situación que ocurrió con gran frecuencia en aquella época. Vende el departamento al adquirente y le propone alquilárselo.

\textbf{Sin el constituto posesorio}, el señor debería entregar la posesión material al adquirente y luego reingresar como inquilino. Sería otro disparate jurídico.

\textbf{Solución legal:} quien era propietario (poseedor a título de dueño) pasa a ser inquilino (tenedor). Desciende de categoría: la flecha va hacia abajo ($\downarrow$).
\end{example}

\section{Cuadro comparativo}

\begin{table}[H]
\centering
\small
\begin{tabularx}{\textwidth}{>{\raggedright\arraybackslash}p{0.2\textwidth}>{\raggedright\arraybackslash}X>{\raggedright\arraybackslash}X}
\toprule
 & \textbf{Traditio brevi manu} & \textbf{Constituto posesorio} \\
\midrule
Artículos & 1892 y 1933 CCyCN & 1892 y 1923 CCyCN \\[0.4em]
Sujeto que cambia de calidad & Quien tenía la cosa como tenedor (inquilino) pasa a poseerla como propietario & Quien poseía como propietario pasa a tener la cosa como tenedor (inquilino) \\[0.4em]
Movimiento de categoría & Ascenso ($\uparrow$) & Descenso ($\downarrow$) \\[0.4em]
Ejemplo & María, la inquilina que compra su departamento & El señor de 2001 que vende y queda como inquilino \\[0.4em]
Tradición material & No requerida & No requerida \\[0.4em]
Constancia notarial & Sí, requerida en la escritura & Sí, requerida con fecha cierta \\
\bottomrule
\end{tabularx}
\caption{Traditio brevi manu y constituto posesorio}
\end{table}

\section{Requisitos de la tradición traslativa de dominio}

No toda entrega material de una cosa implica transferencia de dominio. Para que la tradición sea traslativa de dominio deben concurrir ciertos requisitos, cuya verificación es clave en el análisis de casos concretos.

\begin{table}[H]
\centering
\small
\begin{tabularx}{\textwidth}{>{\raggedright\arraybackslash}p{0.24\textwidth}>{\raggedright\arraybackslash}X>{\raggedright\arraybackslash}p{0.22\textwidth}}
\toprule
\textbf{Requisito} & \textbf{Descripción} & \textbf{Fundamento normativo} \\
\midrule
Legitimación del transmitente &
Quien transmite debe ser propietario. Solo el propietario puede transmitir la propiedad. &
Art.~399 --- \textit{nemo plus iuris} \\[0.5em]
Acto jurídico causal (título suficiente) &
Debe existir una causa válida: compraventa, donación, dación en pago, permuta, etc. &
Art.~1892 CCyCN \\[0.5em]
Capacidad de ambas partes &
Tanto el que transmite como el que adquiere deben ser capaces. &
Art.~1892 CCyCN \\[0.5em]
Escritura pública (inmuebles) &
Si se trata de inmuebles, el título suficiente debe plasmarse en escritura pública. &
Art.~1017 CCyCN \\[0.5em]
Tradición material efectiva &
La entrega física de la cosa, salvo los dos casos de excepción (traditio brevi manu y constituto posesorio). &
Art.~1924 CCyCN \\
\bottomrule
\end{tabularx}
\caption{Requisitos de la tradición traslativa de dominio}
\end{table}

\section{Cuadro Cornell: casos especiales de tradición}

\begin{longtable}{>{\raggedright\arraybackslash}p{0.27\textwidth}>{\raggedright\arraybackslash}p{0.65\textwidth}}
\toprule
\textbf{Preguntas / palabras clave} & \textbf{Notas / respuestas} \\
\midrule
\endfirsthead
\toprule
\textbf{Preguntas / palabras clave} & \textbf{Notas / respuestas} \\
\midrule
\endhead
\bottomrule
\endfoot
\textbf{¿Qué es la traditio brevi manu?} &
Es la excepción que permite que quien ya tenía la cosa como tenedor (ej.: inquilino) pase a tenerla como propietario sin tradición material, con constancia en la escritura. Arts.~1892 y 1933 CCyCN. El vínculo asciende de categoría ($\uparrow$). \\[0.8em]

\textbf{¿Qué es el constituto posesorio?} &
Es el caso inverso: quien era propietario (poseedor) pasa a tener la cosa como tenedor (ej.: inquilino) sin tradición material, con constancia notarial. Arts.~1892 y 1923 CCyCN. El vínculo desciende de categoría ($\downarrow$). \\[0.8em]

\textbf{¿Qué requisitos debe reunir la tradición traslativa de dominio?} &
Legitimación del transmitente (propietario, Art.~399); acto jurídico causal o título suficiente (Art.~1892); capacidad de ambas partes (Art.~1892); escritura pública si se trata de inmuebles (Art.~1017); y tradición material efectiva, salvo los dos casos de excepción (Art.~1924). \\
\midrule
\multicolumn{2}{p{\dimexpr0.92\textwidth+2\tabcolsep\relax}}{%
\textbf{Resumen:} la tradición no puede suplirse por fórmulas escritas, salvo en la traditio brevi manu (el tenedor pasa a ser propietario) y en el constituto posesorio (el propietario pasa a ser tenedor), que no requieren tradición material pero sí constancia notarial. La tradición traslativa de dominio exige además propietario, título suficiente, capacidad de las partes y escritura pública en inmuebles.} \\
\end{longtable}

% ==========================================================
% CAPÍTULO 4
% ==========================================================
\chapter[Publicidad, adquisición legal, transmisión y extinción]{Publicidad, adquisición legal originaria, transmisión y extinción}

\section{Publicidad de los derechos reales}

Los derechos reales deben ser respetados por todos y, por tanto, deben ser \textbf{conocidos por todos}.

\begin{important}
\textbf{Artículo 1893.} Establece el principio de publicidad de los derechos reales: para que sean oponibles a terceros deben ser conocidos. La publicidad puede ser posesoria (tener la cosa) o registral (inscripción en el registro público).
\end{important}

\begin{table}[H]
\centering
\small
\begin{tabularx}{\textwidth}{>{\raggedright\arraybackslash}p{0.25\textwidth}>{\raggedright\arraybackslash}X}
\toprule
\textbf{Tipo de publicidad} & \textbf{Descripción} \\
\midrule
Publicidad posesoria &
Desde la existencia misma del hombre en sociedad, el vínculo más claro con una cosa es tenerla. La posesión es en sí misma publicidad. \\[0.5em]
Publicidad registral &
La inscripción en un registro público hace oponible el derecho real frente a terceros. Sus efectos varían: declarativo en inmuebles, constitutivo en automotores. \\
\bottomrule
\end{tabularx}
\caption{Tipos de publicidad de los derechos reales}
\end{table}

\section{Adquisición legal originaria}

El nuevo Código contempla casos de adquisición originaria en los que el derecho real se adquiere directamente por disposición legal, sin que concurran título y modo en el sentido tradicional. Estos casos presentan especial interés por su vinculación con el orden público y el derecho de familia.

\textbf{Artículo 1894.} Regula los casos de adquisición legal originaria de derechos reales: el derecho se adquiere porque la ley lo dispone, sin necesidad de título y modo en sentido estricto.

Los casos de adquisición legal originaria mencionados son:
\begin{enumerate}
    \item el derecho real de habitación del cónyuge supérstite;
    \item el derecho real de habitación del conviviente supérstite;
    \item los condominios con indivisión forzosa;
    \item los derechos de adquirentes de buena fe con título oneroso (cuyo desarrollo corresponde a clases posteriores).
\end{enumerate}

\begin{important}
El reconocimiento del conviviente como titular de un derecho real de habitación es uno de los cambios de paradigma más importantes del Código Civil y Comercial respecto del régimen anterior.
\end{important}

\section{Transmisión de los derechos reales}

Todos los derechos, reales o personales, son transmisibles como regla general. La transmisión de derechos reales se rige por la teoría del título y modo: quien transmite lo hace con título suficiente y modo, y quien adquiere también los recibe de ese modo. Los casos particulares se analizan al tratar cada derecho real en especial.

\section{Extinción de los derechos reales}

\textbf{Artículo 1907.} El derecho real se extingue de dos formas: relativa y absoluta.

\begin{table}[H]
\centering
\small
\begin{tabularx}{\textwidth}{>{\raggedright\arraybackslash}p{0.2\textwidth}>{\raggedright\arraybackslash}X>{\raggedright\arraybackslash}X}
\toprule
\textbf{Tipo de extinción} & \textbf{Descripción} & \textbf{Ejemplo} \\
\midrule
Extinción relativa &
El derecho real se extingue para un titular, pero continúa en cabeza de otro. &
Venta, donación, dación en pago. El dominio continúa, pero en otra persona. \\[0.5em]
Extinción absoluta &
El derecho real se extingue en forma definitiva, sin que subsista en ningún sujeto. &
Destrucción de la cosa. Al desaparecer el objeto del vínculo, el derecho real cesa definitivamente. \\
\bottomrule
\end{tabularx}
\caption{Extinción relativa y absoluta (Art.~1907)}
\end{table}

También se menciona la \textbf{consolidación} como otra forma de extinción: ocurre cuando una persona que tenía la cosa en virtud de otro derecho real (por ejemplo, el usufructuario) adquiere la propiedad de ella, reuniendo en una misma persona la titularidad de ambos derechos, de modo que el derecho menor queda absorbido.

\section{El derecho real como relación de poder}

El derecho real es una relación de poder del titular respecto de la cosa. Ese es el esquema, la línea o flecha que permite perseguirla en ocasiones, como ocurre con las hipotecas, aun cuando hubiera cambiado de titular. Esa relación de poder es defendida y protegida por el sistema jurídico y debe ser respetada por toda la sociedad. De allí la indefectible necesidad de la publicidad, ya sea posesoria o registral.

Los seres humanos tienen, desde el nacimiento y a lo largo de su existencia, contactos permanentes con las cosas. Ese contacto directo e inmediato tiene una protección jurídica que debe ser respetada, porque es también una manera de respetar la paz social.

No siempre los vínculos de poder están en cabeza del titular del derecho real: no siempre las personas viven en las tierras de las que son propietarias. Existen diversas relaciones de poder y diversas calidades en que un sujeto puede vincularse con una cosa.

\begin{note}
Como contenidos posteriores se anuncian el estudio de todas las relaciones de poder con las cosas (posesión, tenencia y sus variantes) y la profundización de los casos de adquisición legal originaria y de los supuestos de buena fe con título oneroso.
\end{note}

\section{Cuadro Cornell: publicidad, adquisición legal y extinción}

\begin{longtable}{>{\raggedright\arraybackslash}p{0.27\textwidth}>{\raggedright\arraybackslash}p{0.65\textwidth}}
\toprule
\textbf{Preguntas / palabras clave} & \textbf{Notas / respuestas} \\
\midrule
\endfirsthead
\toprule
\textbf{Preguntas / palabras clave} & \textbf{Notas / respuestas} \\
\midrule
\endhead
\bottomrule
\endfoot
\textbf{¿Qué establece el Art.~1893 sobre la publicidad?} &
Los derechos reales, para ser oponibles a terceros, deben ser conocidos. La publicidad puede ser posesoria (tener la cosa) o registral (inscripción en el registro público). \\[0.8em]

\textbf{¿Qué es la adquisición legal originaria (Art.~1894)?} &
Son casos en que la ley dispone directamente la adquisición del derecho real, sin requerir título y modo. Ejemplos: derecho real de habitación del cónyuge supérstite y del conviviente supérstite, condominios con indivisión forzosa. \\[0.8em]

\textbf{¿Cómo se extingue el derecho real?} &
De dos formas (Art.~1907): (1) relativa: el derecho se extingue para el titular pero continúa en otro (ej.: venta); (2) absoluta: el derecho se extingue definitivamente, generalmente por destrucción de la cosa. También por consolidación (reunión en un solo titular). \\
\midrule
\multicolumn{2}{p{\dimexpr0.92\textwidth+2\tabcolsep\relax}}{%
\textbf{Resumen:} los derechos reales se hacen oponibles mediante la publicidad posesoria o registral (Art.~1893). Existen casos de adquisición legal originaria sin título y modo (Art.~1894). La extinción puede ser relativa o absoluta (Art.~1907), y también opera por consolidación.} \\
\end{longtable}

% ==========================================================
% SÍNTESIS FINAL
% ==========================================================
\chapter*{Síntesis final}
\addcontentsline{toc}{chapter}{Síntesis final}
\markboth{Síntesis final}{Síntesis final}

Los derechos reales constituyen vínculos jurídicos directos e inmediatos entre un titular y una cosa, regulados exclusivamente por la ley federal bajo normas de orden público. El Código Civil y Comercial de la Nación, imbuido de la constitucionalización e internacionalización del derecho privado, reconoce exactamente 14 derechos reales y establece el sistema de su adquisición, transmisión y extinción.

La adquisición derivada por actos entre vivos exige la concurrencia de \textbf{título suficiente} (el acto jurídico causal: compraventa, donación, dación en pago) y \textbf{modo suficiente} (la tradición: entrega material de la cosa), conforme al artículo 1892, sin que la tradición pueda suplirse por fórmulas escritas, salvo los dos casos de excepción legalmente previstos: la \textbf{traditio brevi manu} (el tenedor que pasa a ser propietario) y el \textbf{constituto posesorio} (el propietario que pasa a ser tenedor).

Para su oponibilidad a terceros, los derechos reales se publicitan mediante la posesión o la inscripción registral, cuyo efecto es declarativo en inmuebles y constitutivo en automotores. El sistema se cierra con la extinción del derecho, que puede ser relativa (continúa en otro titular) o absoluta (se extingue definitivamente, típicamente por destrucción de la cosa), siempre en el marco de la preeminencia del interés colectivo sobre el individual que proclama el artículo 240 del Código.

\end{document}
